% =========================================================================
% UNIT 4: PROBABILITY, RANDOM VARIABLES, AND DISTRIBUTIONS
% =========================================================================
\chapter{Unit 4: Probability, Random Variables, \& Probability Distributions}

\section{Unit Overview \& Exam Weight}
Unit 4 accounts for \textbf{10\% to 20\%} of the AP Statistics Exam. Probability theory provides the vital mathematical bridge connecting descriptive data exploration (Units 1--3) with statistical inference (Units 6--9). Mastery of random variable transformations, independence, Binomial models, and Geometric distributions is essential for scoring a 5.

\subsection*{Key Unit Vocabulary}
\begin{description}[leftmargin=1.5em, style=nextline]
  \item[Law of Large Numbers:] As the number of repetitions of a chance experiment increases, the relative frequency of an event approaches its true theoretical probability.
  \item[Sample Space ($S$):] The set of all possible outcomes of a random phenomenon.
  \item[Disjoint (Mutually Exclusive) Events:] Events that cannot occur simultaneously: $P(A \cap B) = 0$.
  \item[Independent Events:] Two events $A$ and $B$ are independent if knowing that one event has occurred does not change the probability that the other event occurs: $P(A \mid B) = P(A)$ and $P(B \mid A) = P(B)$.
  \item[Conditional Probability:] The probability that one event happens given that another event is already known to have occurred: $P(A \mid B) = P(A \cap B) / P(B)$.
  \item[Discrete Random Variable ($X$):] Takes a fixed set of countable values with gaps between them, defined by a probability distribution where $\sum P(x_i) = 1$.
  \item[Expected Value ($E(X) = \mu_X$):] The long-run average outcome of a random variable over many repeated trials: $\mu_X = \sum x_i P(x_i)$.
  \item[Variance ($\text{Var}(X) = \sigma_X^2$):] The average squared deviation from the mean: $\sigma_X^2 = \sum (x_i - \mu_X)^2 P(x_i)$.
  \item[Binomial Setting ($B-I-N-S$):] Binary outcomes, Independent trials, Number of trials $n$ fixed, Same probability of success $p$ for each trial.
  \item[Geometric Setting ($B-I-T-S$):] Binary outcomes, Independent trials, Trials continue until the first success, Same probability of success $p$.
\end{description}

\section{Core Concept Lessons}

\begin{lessonbox}[Lesson 4.1: Fundamental Probability Axioms & The Independence Test]
\begin{itemize}[leftmargin=1.2em]
  \item \textbf{Complement Rule:} $P(A^c) = 1 - P(A)$.
  \item \textbf{General Addition Rule:} $P(A \cup B) = P(A) + P(B) - P(A \cap B)$.
  \item \textbf{Conditional Probability Formula:}
  \begin{equation}
    P(A \mid B) = \frac{P(A \cap B)}{P(B)}, \qquad P(A \cap B) = P(B) \cdot P(A \mid B)
  \end{equation}
  \item \textbf{The Official AP Independence Verification Test:}
  Two events $A$ and $B$ are independent if and only if any of the following equivalence conditions hold:
  \begin{equation}
    P(A \mid B) = P(A) \quad \text{OR} \quad P(B \mid A) = P(B) \quad \text{OR} \quad P(A \cap B) = P(A) \cdot P(B)
  \end{equation}
\end{itemize}
\end{lessonbox}

\begin{trapbox}[Exam Trap Alert: Disjoint vs. Independent Events]
Disjoint (mutually exclusive) events and Independent events are \textbf{NEVER} the same thing! In fact, two events with non-zero probabilities can \textbf{never} be both disjoint and independent!
\begin{itemize}[leftmargin=1.2em]
  \item If $A$ and $B$ are \textbf{disjoint}, they cannot happen together ($P(A \cap B) = 0$). If $A$ happens, the probability that $B$ happens drops to exactly 0! Thus, knowing $A$ occurred completely changes the probability of $B$, proving disjoint events are \textbf{heavily dependent}!
  \item On FRQs, never write \textit{"Events A and B are independent because they cannot happen at the same time."} That earns an immediate \textbf{Incomplete (I)} score!
\end{itemize}
\end{trapbox}

\begin{lessonbox}[Lesson 4.2: Discrete Random Variables & Linear Transformations]
\begin{equation}
  \mu_X = E(X) = \sum_{i=1}^k x_i P(x_i), \qquad \sigma_X^2 = \sum_{i=1}^k (x_i - \mu_X)^2 P(x_i), \qquad \sigma_X = \sqrt{\sigma_X^2}
\end{equation}
\textbf{Linear Transformation of a Single Random Variable ($Y = a + bX$):}
\begin{equation}
  \mu_{a + bX} = a + b\mu_X, \qquad \sigma_{a + bX}^2 = b^2 \sigma_X^2, \qquad \sigma_{a + bX} = |b|\sigma_X
\end{equation}
Adding a constant shifts the center but leaves spread completely untouched.
\end{lessonbox}

\begin{lessonbox}[Lesson 4.3: Linear Combinations of Independent Random Variables]
Let $X$ and $Y$ be two \textbf{independent} random variables:
\begin{equation}
  \mu_{X + Y} = \mu_X + \mu_Y, \qquad \mu_{X - Y} = \mu_X - \mu_Y
\end{equation}
\begin{equation}
  \sigma_{X + Y}^2 = \sigma_X^2 + \sigma_Y^2, \qquad \sigma_{X - Y}^2 = \sigma_X^2 + \sigma_Y^2
\end{equation}
\begin{equation}
  \sigma_{X \pm Y} = \sqrt{\sigma_X^2 + \sigma_Y^2}
\end{equation}
\textbf{The Universal Rule of Variances:} \textit{"Standard deviations NEVER add directly! Variances ALWAYS add, even when subtracting two random variables!"}
\end{lessonbox}

\begin{lessonbox}[Lesson 4.4: The Binomial Probability Distribution (B-I-N-S)]
A random variable $X \sim B(n, p)$ represents the number of successes in $n$ independent trials:
\begin{enumerate}[leftmargin=1.5em]
  \item \textbf{B --- Binary:} Each trial has exactly two outcomes: Success ($p$) or Failure ($1-p$).
  \item \textbf{I --- Independent:} Trials are mutually independent.
  \item \textbf{N --- Number:} The number of trials $n$ is fixed in advance.
  \item \textbf{S --- Same Probability:} Probability of success $p$ is constant for each trial.
\end{enumerate}
\begin{equation}
  P(X = k) = \binom{n}{k} p^k (1-p)^{n-k} = \frac{n!}{k!(n-k)!} p^k (1-p)^{n-k}
\end{equation}
\begin{equation}
  \mu_X = np, \qquad \sigma_X = \sqrt{np(1-p)}
\end{equation}
\textbf{10\% Condition for Binomial Sampling:} When sampling without replacement from a finite population of size $N$, trials are approximately independent if $n \le 0.10N$.
\end{lessonbox}

\begin{lessonbox}[Lesson 4.5: The Geometric Probability Distribution (B-I-T-S)]
A random variable $Y \sim G(p)$ counts the number of trials required until the \textbf{first success} occurs:
\begin{equation}
  P(Y = k) = (1-p)^{k-1} p
\end{equation}
\begin{equation}
  \mu_Y = E(Y) = \frac{1}{p}, \qquad \sigma_Y = \frac{\sqrt{1-p}}{p}
\end{equation}
\begin{equation}
  P(Y > k) = (1-p)^k \quad \text{(the probability of failing on the first $k$ consecutive trials)}
\end{equation}
\end{lessonbox}

\section{Worked Step-by-Step Examples}

\subsection*{Example 4.1: Combining Independent Random Variables}
\textbf{Problem:} A packaging company ships boxes containing a glass vase and protective packing peanuts.
\begin{itemize}[leftmargin=1.2em]
  \item The weight of the empty shipping box follows $X \sim N(12\text{ oz}, 1.5\text{ oz})$.
  \item The weight of the glass vase follows $Y \sim N(28\text{ oz}, 2.0\text{ oz})$.
  \item The weight of the packing peanuts follows $W \sim N(4\text{ oz}, 0.5\text{ oz})$.
\end{itemize}
Assuming the weights are independent, what is the probability that a randomly assembled total package exceeds $48\text{ oz}$?

\textbf{Solution:}
Let $T = X + Y + W = \text{total package weight}$.
Because the linear combination of independent normal random variables is itself normal:
\[ \mu_T = \mu_X + \mu_Y + \mu_W = 12 + 28 + 4 = 44.0\text{ oz} \]
\[ \sigma_T^2 = \sigma_X^2 + \sigma_Y^2 + \sigma_W^2 = 1.5^2 + 2.0^2 + 0.5^2 = 2.25 + 4.00 + 0.25 = 6.50\text{ oz}^2 \]
\[ \sigma_T = \sqrt{6.50} \approx 2.55\text{ oz} \]
Thus, $T \sim N(44.0, 2.55)$. We seek $P(T > 48)$:
\[ z = \frac{48 - 44.0}{2.55} = \frac{4.0}{2.55} \approx 1.57 \]
$P(T > 48) = P(Z > 1.57) = 1 - 0.9418 = \mathbf{0.0582}$ (approximately \textbf{5.82\%}).

\section{TI-84 Plus CE Calculator Playbook}

\begin{calculatorbox}[TI-84 Playbook: Binomial and Geometric Distributions]
\begin{itemize}[leftmargin=1.2em]
  \item \textbf{Single Binomial Probability $P(X = k)$:} Press \texttt{[2nd] [VARS]} (\texttt{DISTR}) $\rightarrow$ Select \texttt{A:binompdf(}. Enter $n, p, k$.
  \item \textbf{Cumulative Binomial Probability $P(X \le k)$:} Press \texttt{[2nd] [VARS]} $\rightarrow$ Select \texttt{B:binomcdf(}. Enter $n, p, k$.
  \item \textbf{Geometric Probability $P(Y = k)$:} Press \texttt{[2nd] [VARS]} $\rightarrow$ Select \texttt{E:geometpdf(}. Enter $p, k$.
  \item \textbf{Cumulative Geometric Probability $P(Y \le k)$:} Press \texttt{[2nd] [VARS]} $\rightarrow$ Select \texttt{F:geometcdf(}. Enter $p, k$.
\end{itemize}
\end{calculatorbox}

\section{Comprehensive Practice Problem Set}

\subsection*{Part I: 20 Multiple-Choice Questions}

\begin{enumerate}[leftmargin=1.5em]
  \item If $P(A) = 0.40, P(B) = 0.50,$ and $P(A \cap B) = 0.20$, are events $A$ and $B$ independent?
  \begin{enumerate}[label=(\Alph*)]
    \item Yes, because $P(A) \cdot P(B) = 0.40 \times 0.50 = 0.20 = P(A \cap B)$.
    \item No, because $P(A \cap B) \neq 0$.
    \item Yes, because $P(A) + P(B) \neq 1$.
    \item No, because $P(A) \neq P(B)$.
    \item Cannot be determined.
  \end{enumerate}
  \textbf{Answer: (A).} The multiplication rule $P(A \cap B) = P(A)P(B)$ proves independence.

  \item Two events $C$ and $D$ have $P(C) = 0.6$ and $P(D) = 0.3$. If $C$ and $D$ are mutually exclusive, what is $P(C \cup D)$?
  \begin{enumerate}[label=(\Alph*)]
    \item 0.18 \item 0.90 \item 0.72 \item 0.30 \item 0.00
  \end{enumerate}
  \textbf{Answer: (B).} For disjoint events, $P(C \cup D) = P(C) + P(D) = 0.6 + 0.3 = 0.90$.

  \item A random variable $X$ has $\mu_X = 10, \sigma_X = 3$. If $Y = 4X - 5$, what are the mean and standard deviation of $Y$?
  \begin{enumerate}[label=(\Alph*)]
    \item $\mu_Y = 35, \sigma_Y = 12$ \item $\mu_Y = 35, \sigma_Y = 7$ \item $\mu_Y = 40, \sigma_Y = 12$ \item $\mu_Y = 35, \sigma_Y = 48$ \item $\mu_Y = 40, \sigma_Y = 3$
  \end{enumerate}
  \textbf{Answer: (A).} $\mu_Y = 4(10) - 5 = 35$. $\sigma_Y = |4|(3) = 12$.

  \item Independent variables $X$ and $Y$ have $\sigma_X = 4$ and $\sigma_Y = 3$. What is the standard deviation of $D = X - Y$?
  \begin{enumerate}[label=(\Alph*)]
    \item 1 \item 5 \item 7 \item 25 \item $\sqrt{7}$
  \end{enumerate}
  \textbf{Answer: (B).} Variances add: $\sigma_D = \sqrt{\sigma_X^2 + \sigma_Y^2} = \sqrt{16 + 9} = \sqrt{25} = 5$.

  \item 20\% of cereal boxes contain a prize. A student buys 5 boxes. What is the probability of getting exactly 2 prizes?
  \begin{enumerate}[label=(\Alph*)]
    \item 0.040 \item 0.2048 \item 0.3277 \item 0.0512 \item 0.1638
  \end{enumerate}
  \textbf{Answer: (B).} $P(X = 2) = \binom{5}{2} (0.2)^2 (0.8)^3 = 10(0.04)(0.512) = 0.2048$.

  \item In a binomial distribution with $n = 100$ and $p = 0.30$, what are the mean and standard deviation?
  \begin{enumerate}[label=(\Alph*)]
    \item $\mu = 30, \sigma = 21$ \item $\mu = 30, \sigma = 4.58$ \item $\mu = 30, \sigma = 2.10$ \item $\mu = 50, \sigma = 5.00$ \item $\mu = 100, \sigma = 30$
  \end{enumerate}
  \textbf{Answer: (B).} $\mu = np = 30$. $\sigma = \sqrt{np(1-p)} = \sqrt{100(0.3)(0.7)} = \sqrt{21} \approx 4.58$.

  \item A basketball player makes 75\% of her free throws. What is the probability that her first miss occurs on her $4^{\text{th}}$ attempt?
  \begin{enumerate}[label=(\Alph*)]
    \item $(0.75)^3 (0.25)$ \item $(0.75)(0.25)^3$ \item $\binom{4}{1} (0.75)^3 (0.25)$ \item $(0.25)^4$ \item $(0.75)^4$
  \end{enumerate}
  \textbf{Answer: (A).} Geometric distribution: Success = Miss ($p = 0.25$). Trials $= 4 \implies (0.75)^3(0.25) \approx 0.1055$.

  \item What is the expected number of rolls of a fair 6-sided die until the first 6 appears?
  \begin{enumerate}[label=(\Alph*)]
    \item 1 \item 3 \item 3.5 \item 6 \item 36
  \end{enumerate}
  \textbf{Answer: (D).} For a geometric distribution, $\mu = 1/p = 1/(1/6) = 6$.

  \item A probability distribution for a discrete random variable $X$ requires which two conditions?
  \begin{enumerate}[label=(\Alph*)]
    \item $0 \le P(x_i) \le 1$ and $\sum P(x_i) = 1$
    \item $\mu = 0$ and $\sigma = 1$
    \item $P(x_i)$ is symmetric and bell-shaped
    \item $n \ge 30$ and $np \ge 10$
    \item $X$ takes only positive values
  \end{enumerate}
  \textbf{Answer: (A).} Foundational axioms for valid probability distributions.

  \item If $P(A \mid B) = 0.70$ and $P(B) = 0.50$, what is $P(A \cap B)$?
  \begin{enumerate}[label=(\Alph*)]
    \item 0.20 \item 0.35 \item 1.20 \item 0.714 \item 0.25
  \end{enumerate}
  \textbf{Answer: (B).} $P(A \cap B) = P(B) \cdot P(A \mid B) = 0.50 \times 0.70 = 0.35$.
\end{enumerate}

\begin{enumerate}[resume, leftmargin=1.5em]
  \item Which of the following is NOT a requirement for a binomial distribution?
  \begin{enumerate}[label=(\Alph*)]
    \item Fixed number of trials $n$ \item Exactly two outcomes per trial \item Trials continue until the first success \item Independent trials \item Constant probability of success $p$
  \end{enumerate}
  \textbf{Answer: (C).} Waiting for the first success defines a geometric distribution, not binomial.

  \item $X$ and $Y$ are independent random variables with $\mu_X = 50, \mu_Y = 20$. What is $\mu_{X - 2Y}$?
  \begin{enumerate}[label=(\Alph*)]
    \item 10 \item 30 \item 50 \item 90 \item $-10$
  \end{enumerate}
  \textbf{Answer: (A).} $\mu = 50 - 2(20) = 10$.

  \item In problem 12, if $\sigma_X = 5$ and $\sigma_Y = 3$, what is the standard deviation of $X - 2Y$?
  \begin{enumerate}[label=(\Alph*)]
    \item $-1$ \item $\sqrt{25 + 36} = \sqrt{61} \approx 7.81$ \item $\sqrt{25 - 36}$ \item 11 \item 16
  \end{enumerate}
  \textbf{Answer: (B).} $\sigma^2 = \sigma_X^2 + 2^2 \sigma_Y^2 = 25 + 4(9) = 25 + 36 = 61 \implies \sigma = \sqrt{61}$.

  \item The probability of winning a carnival game is $p = 0.05$. What is the probability that you fail on the first 10 plays?
  \begin{enumerate}[label=(\Alph*)]
    \item $(0.05)^{10}$ \item $(0.95)^{10}$ \item $1 - (0.95)^{10}$ \item $10(0.05)$ \item $(0.95)(0.05)^9$
  \end{enumerate}
  \textbf{Answer: (B).} $P(Y > 10) = (1 - p)^{10} = (0.95)^{10} \approx 0.5987$.

  \item Two fair coins are tossed. Let $X = \text{number of heads}$. What is $\text{Var}(X)$?
  \begin{enumerate}[label=(\Alph*)]
    \item 0.25 \item 0.50 \item 0.75 \item 1.00 \item 2.00
  \end{enumerate}
  \textbf{Answer: (B).} Binomial with $n=2, p=0.5 \implies \sigma^2 = np(1-p) = 2(0.5)(0.5) = 0.50$.

  \item If $P(A) = 0.7, P(B) = 0.4,$ and $P(A \cup B) = 0.85$, what is $P(A \cap B)$?
  \begin{enumerate}[label=(\Alph*)]
    \item 0.15 \item 0.25 \item 0.28 \item 0.55 \item 0.00
  \end{enumerate}
  \textbf{Answer: (B).} $P(A \cap B) = P(A) + P(B) - P(A \cup B) = 0.7 + 0.4 - 0.85 = 0.25$.

  \item When sampling from a population of size $N = 500$, what is the maximum sample size $n$ allowed to use the binomial formula safely without replacement?
  \begin{enumerate}[label=(\Alph*)]
    \item 10 \item 30 \item 50 \item 100 \item 250
  \end{enumerate}
  \textbf{Answer: (C).} By the 10\% condition, $n \le 0.10(500) = 50$.

  \item Which of the following statements about independent events $A$ and $B$ is always true?
  \begin{enumerate}[label=(\Alph*)]
    \item $P(A \cup B) = 1$ \item $P(A \mid B) = P(A)$ \item $P(A \cap B) = 0$ \item $P(A) = P(B)$ \item $P(A) + P(B) = 1$
  \end{enumerate}
  \textbf{Answer: (B).} Definition of statistical independence.

  \item A discrete random variable $X$ takes values 1, 2, 3 with probabilities 0.2, 0.5, 0.3. What is $E(X)$?
  \begin{enumerate}[label=(\Alph*)]
    \item 2.0 \item 2.1 \item 2.3 \item 2.5 \item 3.0
  \end{enumerate}
  \textbf{Answer: (B).} $E(X) = 1(0.2) + 2(0.5) + 3(0.3) = 0.2 + 1.0 + 0.9 = 2.1$.

  \item A binomial distribution with $n = 50$ and $p = 0.40$ can be approximated by a normal distribution because:
  \begin{enumerate}[label=(\Alph*)]
    \item $n \ge 30$
    \item $np = 20 \ge 10$ and $n(1-p) = 30 \ge 10$
    \item $p = 0.40 < 0.50$
    \item $n \le 0.10N$
    \item The trials are dependent
  \end{enumerate}
  \textbf{Answer: (B).} The Large Counts condition ($np \ge 10$ and $n(1-p) \ge 10$) justifies normality.
\end{enumerate}

\subsection*{Part II: Free-Response Practice Problem}

\begin{workbookbox}[FRQ 4.1: Binomial Probability & Random Variable Combinations]
\textbf{Problem:} An airline knows that 5\% of passengers who purchase tickets do not show up for their flight. A flight operates on an aircraft with 100 seats, but the airline sells 102 tickets.
(a) Define the random variable and check whether the conditions for a binomial distribution are met.\\
(b) Calculate the probability that more passengers show up than there are seats available.\\
(c) The airline charges \$200 per ticket sold. If more than 100 passengers show up, each bumped passenger is paid \$500 in compensation. What is the expected net revenue for this flight?

\textbf{Grader-Approved Score-4 Solution:}
\begin{itemize}[leftmargin=1.2em]
  \item \textbf{Part (a):} Let $X = \text{number of passengers who show up}$. $X \sim B(102, 0.95)$.\\
  \textit{Binary:} Passenger shows up (Success, $p=0.95$) or does not show up (Failure, $1-p=0.05$).\\
  \textit{Independent:} Passenger decisions are assumed independent.\\
  \textit{Number:} Fixed $n = 102$ tickets.\\
  \textit{Same Probability:} $p = 0.95$ constant for each passenger.
  \item \textbf{Part (b):} More passengers than seats means $X = 101$ or $X = 102$.\\
  $P(X = 101) = \binom{102}{101} (0.95)^{101} (0.05)^1 = 102(0.00569)(0.05) \approx 0.0290$.\\
  $P(X = 102) = \binom{102}{102} (0.95)^{102} (0.05)^0 = (0.95)^{102} \approx 0.0054$.\\
  $P(X > 100) = 0.0290 + 0.0054 = \mathbf{0.0344}$ (approx. \textbf{3.44\%}).
  \item \textbf{Part (c):} Guaranteed revenue from 102 tickets $= 102 \times \$200 = \$20,400$.\\
  If 101 show up (prob 0.0290), 1 passenger bumped $\implies \text{Cost} = \$500$.\\
  If 102 show up (prob 0.0054), 2 passengers bumped $\implies \text{Cost} = \$1,000$.\\
  Expected compensation cost $= 0.0290(\$500) + 0.0054(\$1,000) = \$14.50 + \$5.40 = \$19.90$.\\
  $\text{Expected Net Revenue} = \$20,400 - \$19.90 = \mathbf{\$20,380.10}$.
\end{itemize}
\end{workbookbox}


\begin{frqrubricbox}[FRQ 4.2: Linear Combinations of Random Variables & Binomial Distributions]
\textbf{Problem:} A specialty roaster fills coffee bags using two automated dispensers.
Dispenser 1 dispense volume ($X$) follows $N(\mu_X = 16.2\text{ oz}, \sigma_X = 0.5\text{ oz})$.
Dispenser 2 dispense volume ($Y$) follows $N(\mu_Y = 16.2\text{ oz}, \sigma_Y = 0.5\text{ oz})$.
Assume the dispensers operate independently.
\begin{enumerate}[label=(\alph*)]
  \item Find the mean and standard deviation of the total coffee volume dispensed when a two-bag combo is filled ($T = X + Y$).
  \item Find the probability that the absolute difference between the two bags $|X - Y|$ exceeds 1.0 ounce.
  \item An automated sealing line inspects filled bags. The probability that any single bag has a seal defect is $p = 0.04$. In an independent random batch of 50 bags, calculate the probability that at least 2 bags have seal defects.
\end{enumerate}

\textbf{Grader-Approved Score-4 Solution \& Rubric:}
\begin{itemize}[leftmargin=1.2em]
  \item \textbf{Part (a):} $\mu_T = \mu_X + \mu_Y = 16.2 + 16.2 = 32.4\text{ oz}$.\\
  Because $X$ and $Y$ are independent, variances add:
  $\sigma_T = \sqrt{\sigma_X^2 + \sigma_Y^2} = \sqrt{0.5^2 + 0.5^2} = \sqrt{0.25 + 0.25} = \sqrt{0.50} \approx 0.707\text{ oz}$.
  \item \textbf{Part (b):} Let $D = X - Y$. $\mu_D = 16.2 - 16.2 = 0\text{ oz}$.\\
  $\sigma_D = \sqrt{\sigma_X^2 + \sigma_Y^2} = \sqrt{0.5^2 + 0.5^2} \approx 0.707\text{ oz}$. (Variances ALWAYS add!)\\
  $D \sim N(0, 0.707)$.
  $z = \frac{1.0 - 0}{0.707} \approx 1.414$.\\
  $P(|D| > 1.0) = P(Z > 1.414) + P(Z < -1.414) = 2 \times 0.0786 = \mathbf{0.1573}$.
  \item \textbf{Part (c):} Let $W$ be the number of defective bags. $W \sim B(n = 50, p = 0.04)$.\\
  $P(W \ge 2) = 1 - [P(W = 0) + P(W = 1)]$.\\
  $P(W = 0) = (0.96)^{50} \approx 0.1299$.\\
  $P(W = 1) = \binom{50}{1} (0.04)^1 (0.96)^{49} = 50(0.04)(0.1353) \approx 0.2706$.\\
  $P(W \ge 2) = 1 - (0.1299 + 0.2706) = 1 - 0.4005 = \mathbf{0.5995}$.
\end{itemize}
\end{frqrubricbox}

\begin{trapbox}[Unit 4 Top 5 AP Exam Pitfalls to Avoid]
\begin{enumerate}[leftmargin=1.2em]
  \item \textbf{Subtracting Variances:} Never subtract variances! When subtracting random variables ($X - Y$), variances still add: $\sigma_{X-Y}^2 = \sigma_X^2 + \sigma_Y^2$.
  \item \textbf{Confusing Binomial with Geometric:} Binomial counts successes in fixed $n$ trials; Geometric counts trials until first success.
  \item \textbf{Forgetting to State Independence Before Multiplying:} $P(A \cap B) = P(A)P(B)$ only applies when events are independent.
  \item \textbf{Confusing Disjoint (Mutually Exclusive) with Independent:} Disjoint events can NEVER be independent (if one occurs, the other cannot!).
  \item \textbf{Multiplying Standard Deviations by Negative Numbers:} Linear transformation $\sigma_{a+bx} = |b|\sigma_x$. Standard deviations are always non-negative.
\end{enumerate}
\end{trapbox}
\section{Unit 4 Master Summary}
\begin{lessonbox}[Unit 4 Essential Review Checklist]
\begin{itemize}[leftmargin=1.2em]
  \item \textbf{Independence Test:} $P(A \mid B) = P(A)$ or $P(A \cap B) = P(A)P(B)$.
  \item \textbf{Combining RVs:} $\mu_{X \pm Y} = \mu_X \pm \mu_Y$; $\sigma_{X \pm Y}^2 = \sigma_X^2 + \sigma_Y^2$ (variances always add!).
  \item \textbf{Binomial ($BINS$):} $P(X=k) = \binom{n}{k}p^k(1-p)^{n-k}$, $\mu=np$, $\sigma=\sqrt{np(1-p)}$.
  \item \textbf{Geometric ($BITS$):} $P(Y=k) = (1-p)^{k-1}p$, $\mu=1/p$.
\end{itemize}
\end{lessonbox}
\clearpage
